\documentclass[a4paper,12pt]{article}

\usepackage{alltt, fancyvrb, url}
\usepackage{graphicx}
\usepackage[utf8]{inputenc}
\usepackage{float}
\usepackage{hyperref}
\usepackage[italian]{babel}
\usepackage[table]{xcolor}
\usepackage{array}
\usepackage{multirow}
\usepackage{titlesec}
\usepackage{listings}
\usepackage[italian]{cleveref}
\usepackage{acronym}
\usepackage{subcaption}

\renewcommand{\thesection}{\arabic{section}}

\sloppy
\setlength{\parindent}{0pt}

\title{\textbf{Stime}}
\author{}
\date{}

\begin{document}
\maketitle
\tableofcontents

\newpage

\section{Stima della quantità di lavoro}
Per ogni task della WBS è stata effettuata una stima della quantità di lavoro necessaria per completarla, utilizzando
una tecnica \textit{consensus-based}. Nello specifico, per i task con una bassa complessità e grado di incertezza è
stato utilizzato il metodo \textit{Three-Point}, basandosi sulle tre stime (ottimistica, pessimistica e più probabile)
di una figura esperta in quell'attività, per poi calcolare la stima finale come media ponderata delle tre stime,
utilizzando la formula:
\begin{center}
    $E = \frac{O + 4M + P}{6}$
\end{center}
Mentre per i task con una complessità e grado di incertezza più elevati è stato utilizzata la tecnica \textit{Delphi},
nella quale i componenti principali del team forniscono le proprie stime, ed in maniera iterativa si discute per un
massimo di 3 round, fino a convergere ad una stima finale che equivale a quella più votata.

\subsection{Attività 1.2.1}
Stimiamo con il metodo \textit{Three-Point}:
\begin{itemize}
    \item Stima ottimistica (O): 2 giorni
    \item Stima più probabile (M): 3 giorni
    \item Stima pessimistica (P): 5 giorni
\end{itemize}
\begin{center}
    $E = \frac{2 + 4*3 + 5}{6} = 3.17$ giorni
\end{center}

\subsection{Attività 4.2.2}
Stimiamo con la tecnica \textit{Delphi}:

\begin{table}[!htb]
    \footnotesize
    \centering
    
    \begin{subtable}[t]{0.3\textwidth}
        \centering
        \begin{tabular}[t]{|c|c|}
            \hline
            Stima (giorni) & Voti \\
            \hline
            4 & 1          \\
            \hline
            6 & 1             \\
            \hline
            7 & 2             \\
            \hline
            10 & 1              \\
            \hline
            14 & 1             \\    
            \hline
        \end{tabular}
        \caption{Round 1}
    \end{subtable}
    \hfill
    \begin{subtable}[t]{0.3\textwidth}
        \centering
        \begin{tabular}[t]{|c|c|}
            \hline
            Stima (giorni) & Voti \\
            \hline
            4 & 1          \\
            \hline
            6 & 1             \\
            \hline
            7 & 1             \\
            \hline
            8 & 2              \\
            \hline
            10 & 1             \\    
            \hline
        \end{tabular}
        \caption{Round 2}
    \end{subtable}
    \hfill
    \begin{subtable}[t]{0.3\textwidth}
        \centering
        \begin{tabular}[t]{|c|c|}
            \hline
            Stima (giorni) & Voti \\
            \hline
            6 & 2             \\
            \hline
            7 & 3             \\
            \hline
            8 & 1             \\
            \hline
        \end{tabular}
        \caption{Round 3}
    \end{subtable}
\end{table}

Dunque la stima finale è di 7 giorni, in quanto è la stima più votata al termine del terzo round.

\section{Stima delle risorse necessarie}
Il team di sviluppo è composto da 10 membri, di cui 5 sono specializzati nel frontend e 5 nel backend, e i team sono
organizzati in maniera orizzontale. Per questa ragione, i task saranno assegnati durante l'esecuzione del progetto in
base alla disponibilità dei membri del team. Tuttavia, è stata effettuata una stima del numero di risorse (sviluppatori)
necessarie per ogni task. Nella \Cref{table:resources_estimated} sono riportati alcuni esempi di stime.

\begin{table}[!htb]
    \centering
    \begin{tabular}{|c|c|c|}
        \hline
        Attività & Categoria & Numero di risorse \\
        \hline
        1.1.1   & Backend  & 1 Senior + 1 Junior  \\
        \hline
        1.1.2.1 & Backend  & 1 Senior  \\
        \hline
        1.2.1   & Frontend & 1 Junior  \\
        \hline
        4.2.1   & Frontend & 1 Senior + 1 Junior  \\
        \hline
        4.2.2   & Backend  & 1 Senior + 1 Junior  \\
        \hline
        5.1.2   & Backend  & 1 Senior + 2 Junior  \\
        \hline
    \end{tabular}
    \caption{Stima del numero di risorse necessarie per ogni attività.}
    \label{table:resources_estimated}
\end{table}

\section{Stima della durata}
Per stimare la durata di ogni task ci si basa sulla stima della quantità di lavoro effettuata e sul numero di risorse
stimate. Nel caso in cui il numero di risorse stimate sia 1, la durata del task è uguale alla stima della quantità di
lavoro, mentre se il numero di risorse è maggiore, la durata del task consisterebbe nel rapporto tra la stima della
quantità di lavoro e il numero di risorse, ma visto che la relazione non è lineare, si aggiusta in base ai dati storici.
Nella \Cref{table:duration_estimated} sono riportati alcuni esempi di stime.

\begin{table}[!htb]
    \centering
    \begin{tabular}{|c|c|}
        \hline
        Attività & Durata (giorni) \\
        \hline
        1.2.1    & 3  \\
        \hline
        4.2.2    & 5  \\
        \hline
    \end{tabular}
    \caption{Stima della durata di ogni attività.}
    \label{table:duration_estimated}
\end{table}

\section{Stima dei costi}
Per stimare i costi del progetto si è considerato il costo orario di ogni risorsa (50 euro per le risorse senior e di
25 euro per le risorse junior), considerando che la durata delle attività è stata stimata in giorni lavorativi di 8
ore, perché tutti i membri del team sono a tempo pieno sul progetto. \\
Nella \Cref{table:cost_estimated} sono riportati alcuni esempi di stime, in cui il costo di ogni attività è calcolato
come il prodotto tra il costo delle risorse stimate (differenziato per tipo di risorsa) e la durata stimata
dell'attività.

\begin{table}[!htb]
    \centering
    \begin{tabular}{|c|c|}
        \hline
        Attività & Costo (euro) \\
        \hline
        1.2.1    & $25 \times 8 \times 3 = 600$  \\
        \hline
        4.2.2    & $(50 + 25) \times 8 \times 5 = 3000$  \\
        \hline
    \end{tabular}
    \caption{Stima dei costi di ogni attività.}
    \label{table:cost_estimated}
\end{table}

\section{Cash Flow Management}
Dopo aver stimato i costi di tutte le attività, si è stimato che il costo totale del progetto è di 140.000 euro.
Data la richiesta del committente di rimanere entro un budget di 200.000 euro, la nostra proposta economica per la
realizzazione del progetto consiste in 180.000 euro. \\
Il piano di pagamento proposto prevede un anticipo del 20\% alla firma del contratto, e il restante suddiviso in:
\begin{itemize}
    \item 25\% in corrispondenza della milestone 1;
    \item 25\% in corrispondenza della milestone 2;
    \item 30\% alla consegna del progetto.
\end{itemize}

\end{document}